% -*- coding: utf-8 -*-
\chapter{Glosario}

Adaptador: Implementación concreta de un puerto hexagonal. En este trabajo, los adaptadores de salida conectan el RAG con pgvector, el índice BM25, Ollama y el Cross-Encoder; los de entrada son las rutas FastAPI que delegan en los casos de uso.

ANN (Approximate Nearest Neighbors): Técnica de búsqueda aproximada que permite localizar de forma eficiente los vectores más similares en bases de datos de gran tamaño, reduciendo significativamente el tiempo de consulta respecto a una búsqueda exhaustiva.

Arquitectura hexagonal (Ports and Adapters): Estilo arquitectónico propuesto por Cockburn (2005) en el que el dominio de negocio no depende de frameworks ni de infraestructura. El núcleo se comunica con el exterior mediante puertos (interfaces); los adaptadores implementan esos puertos hacia HTTP, bases de datos o modelos de lenguaje.

Backpropagation (Retropropagación): Algoritmo de entrenamiento utilizado en redes neuronales artificiales que ajusta los pesos de las conexiones minimizando el error mediante el cálculo del gradiente de la función de pérdida.

Base de datos vectorial (Vector Database): Sistema especializado para almacenar, indexar y recuperar vectores de alta dimensión generados por modelos de embeddings. Facilita búsquedas semánticas rápidas mediante métricas de similitud.

PostgreSQL/pgvector: Base de datos vectorial de código abierto diseñada para aplicaciones de inteligencia artificial y sistemas RAG. Destaca por su facilidad de integración con Python y LangChain.

Chunk: Fragmento de un documento utilizado como unidad mínima de procesamiento en sistemas RAG. Cada fragmento conserva parte del contexto original para facilitar la recuperación de información.

Chunking: Proceso de dividir documentos extensos en fragmentos más pequeños y manejables antes de generar sus embeddings e indexarlos en una base de datos vectorial.

CLIP (Contrastive Language-Image Pretraining): Modelo multimodal desarrollado por OpenAI capaz de representar imágenes y texto dentro del mismo espacio vectorial, permitiendo búsquedas cruzadas entre ambos formatos.

Context Precision: Métrica utilizada para evaluar la proporción de fragmentos recuperados que realmente son relevantes para responder una consulta.

Context Recall: Métrica que mide la capacidad de un sistema RAG para recuperar toda la información necesaria existente en el corpus documental.

Cosine Similarity (Similitud del coseno): Medida matemática utilizada para calcular el grado de semejanza entre dos vectores mediante el ángulo formado entre ellos.

Cross Encoder: Modelo de reranking que evalúa conjuntamente una consulta y un documento para calcular una puntuación de relevancia más precisa que los modelos bi-encoder

Docker: Plataforma de virtualización ligera basada en contenedores que permite desplegar aplicaciones de forma reproducible e independiente del sistema operativo.

Embedding: Representación numérica de alta dimensión que captura el significado semántico de palabras, frases, documentos o imágenes para permitir comparaciones matemáticas entre ellos.

Faithfulness: Métrica de evaluación que determina el grado en que una respuesta generada por un modelo está fundamentada exclusivamente en la información recuperada del contexto documental.

FAISS (Facebook AI Similarity Search): Biblioteca desarrollada por Meta para realizar búsquedas eficientes de similitud vectorial en colecciones de gran tamaño.

Fine-Tuning: Proceso mediante el cual un modelo de lenguaje previamente entrenado se adapta a una tarea específica utilizando un conjunto adicional de datos especializados.

GPT (Generative Pre-trained Transformer): Familia de modelos de lenguaje basados en la arquitectura Transformer capaces de comprender y generar texto mediante aprendizaje profundo.

GPT-4V (GPT-4 Vision): Variante multimodal de GPT que incorpora capacidades de comprensión tanto de texto como de imágenes.

GraphRAG: Arquitectura RAG que combina recuperación documental con grafos de conocimiento para representar relaciones complejas entre entidades.

Guardrails: Conjunto de reglas, validaciones y mecanismos de control implementados para evitar respuestas incorrectas, peligrosas o alucinadas generadas por modelos de lenguaje.

Hallucination (Alucinación): Fenómeno por el cual un modelo de lenguaje genera información falsa, inventada o no respaldada por los datos disponibles, presentándola como verdadera.

HNSW (Hierarchical Navigable Small World): Algoritmo de indexación utilizado en bases de datos vectoriales que organiza los vectores en grafos jerárquicos para acelerar las búsquedas por similitud.

IA (Inteligencia Artificial): Disciplina de la informática dedicada al desarrollo de sistemas capaces de realizar tareas que normalmente requieren inteligencia humana.

LangChain: Framework de desarrollo que facilita la construcción de aplicaciones basadas en modelos de lenguaje, especialmente arquitecturas RAG y agentes inteligentes.

Layout Analysis: Técnica de visión artificial utilizada para identificar automáticamente títulos, tablas, imágenes y otros elementos estructurales dentro de documentos PDF.

LlamaIndex: Framework especializado en la indexación y recuperación documental para aplicaciones basadas en modelos de lenguaje e inteligencia artificial generativa.

LLM (Large Language Model): Modelo de lenguaje entrenado sobre enormes volúmenes de texto capaz de comprender, generar y razonar utilizando lenguaje natural.

Milvus: Base de datos vectorial de código abierto optimizada para almacenar miles de millones de embeddings con alto rendimiento.

Modelo multimodal: Modelo de inteligencia artificial capaz de procesar simultáneamente distintos tipos de información, como texto, imágenes, tablas, audio o vídeo.

Multi-Head Attention: Mecanismo de la arquitectura Transformer que permite atender simultáneamente diferentes relaciones semánticas presentes dentro de una secuencia.

NLP (Natural Language Processing): Rama de la inteligencia artificial dedicada al procesamiento automático del lenguaje humano.

NLU (Natural Language Understanding): Subdisciplina del NLP centrada en la comprensión del significado, intención y contexto del lenguaje natural.

NLG (Natural Language Generation): Área del NLP encargada de generar automáticamente texto coherente y comprensible a partir de información estructurada.

OCR (Optical Character Recognition): Tecnología que convierte imágenes o documentos escaneados en texto editable mediante reconocimiento óptico de caracteres.

Ollama: Plataforma para ejecutar modelos de lenguaje de forma local, preservando la privacidad y evitando la dependencia de servicios en la nube.

PAC (Política Agraria Común): Política de la Unión Europea destinada a apoyar el sector agrícola mediante ayudas económicas, regulación del mercado y desarrollo rural.

PDF Parsing: Proceso de extracción automática de texto, tablas e imágenes desde documentos PDF para su posterior procesamiento.

Pinecone: Servicio gestionado de base de datos vectorial diseñado para aplicaciones de inteligencia artificial y búsqueda semántica a gran escala.

POSEI (Programa de Opciones Específicas por la Lejanía y la Insularidad): Programa europeo destinado a compensar las dificultades derivadas de la insularidad y ultraperiferia en determinadas regiones agrícolas.

Puerto: Contrato (interfaz) que el dominio o los casos de uso exigen al exterior, sin conocer la tecnología concreta. En el núcleo RAG los puertos salientes cubren embeddings, LLM, repositorio de chunks, índice léxico y reranker.

Prompt Engineering: Disciplina orientada al diseño y optimización de instrucciones dirigidas a modelos de lenguaje con el fin de obtener respuestas más precisas.

PyMuPDF: Biblioteca de Python utilizada para extraer texto, imágenes y metadatos desde documentos PDF.

Qdrant: Base de datos vectorial de código abierto desarrollada en Rust, optimizada para búsquedas semánticas con filtrado avanzado de metadatos.

RAG (Retrieval-Augmented Generation): Arquitectura que combina recuperación de información documental con generación de texto mediante modelos de lenguaje para producir respuestas fundamentadas en fuentes externas.

RAG Multimodal: Evolución del RAG tradicional que incorpora recuperación y procesamiento conjunto de texto, tablas, imágenes, diagramas y otros contenidos multimodales.

RAGAS (Retrieval Augmented Generation Assessment): Framework utilizado para evaluar automáticamente sistemas RAG mediante métricas como fidelidad, precisión contextual y relevancia de las respuestas.

Re-ranking: Etapa posterior a la recuperación documental donde los resultados iniciales se reordenan según su relevancia utilizando modelos más precisos.

Self-Attention (Autoatención): Mecanismo fundamental de los Transformers que permite calcular la importancia relativa entre todos los elementos de una secuencia.

Token: Unidad mínima de información procesada por un modelo de lenguaje, que puede corresponder a una palabra completa, una parte de ella o un signo de puntuación.

Tokenización: Proceso de dividir un texto en tokens antes de ser procesado por un modelo de lenguaje.

Transformer: Arquitectura de redes neuronales presentada en 2017 basada en mecanismos de autoatención, considerada el fundamento de los actuales modelos de lenguaje de gran tamaño.

TruLens: Framework empleado para monitorizar y evaluar el rendimiento, precisión y seguridad de aplicaciones basadas en modelos de lenguaje

Vectorización: Proceso mediante el cual un documento, imagen o texto se transforma en un embedding para su almacenamiento y recuperación semántica.

Vision Transformer (ViT): Arquitectura basada en Transformers adaptada al procesamiento de imágenes mediante la división de estas en pequeños fragmentos o patches.



